\UnnumberedChapter{LỜI MỞ ĐẦU}

Trong bối cảnh công nghệ thông tin ngày càng phát triển, trí tuệ nhân tạo (AI) và học máy (Machine Learning) đang được ứng dụng rộng rãi trong nhiều lĩnh vực, đặc biệt là giáo dục. Việc ứng dụng học máy vào phân tích dữ liệu giúp nâng cao hiệu quả xử lý thông tin, hỗ trợ đánh giá và cải thiện chất lượng giảng dạy, học tập.

Phản hồi của sinh viên là nguồn thông tin quan trọng, phản ánh những trải nghiệm và đánh giá của người học đối với hoạt động giáo dục. Tuy nhiên, việc phân tích thủ công một lượng lớn phản hồi thường mất nhiều thời gian và công sức. Vì vậy, ứng dụng các thuật toán học máy để tự động phân loại cảm xúc từ phản hồi của sinh viên là một hướng nghiên cứu có ý nghĩa thực tiễn.

Xuất phát từ thực tế đó, nhóm chúng em thực hiện đề tài “Phân loại cảm xúc từ phản hồi của sinh viên” trong khuôn khổ học phần Nhập môn Học máy, dưới sự hướng dẫn của ThS. Lê Thị Thùy Trang. Đề tài tập trung vào quá trình khám phá dữ liệu, tiền xử lý văn bản, trích xuất đặc trưng, xây dựng và đánh giá các mô hình phân loại cảm xúc bằng những thuật toán học máy phù hợp.

Thông qua đề tài, nhóm mong muốn vận dụng kiến thức đã học vào thực tiễn, đồng thời nâng cao kỹ năng lập trình, phân tích dữ liệu, xây dựng mô hình và làm việc nhóm. Bên cạnh đó, kết quả nghiên cứu có thể góp phần làm cơ sở tham khảo cho việc ứng dụng học máy trong phân tích phản hồi và nâng cao chất lượng giáo dục.

Do kiến thức và kinh nghiệm còn hạn chế, báo cáo khó tránh khỏi những thiếu sót. Nhóm chúng em rất mong nhận được sự góp ý của cô để tiếp tục hoàn thiện đề tài.